Large language models are increasingly used in settings that require nuanced judgments about emotional tone, reader response, and narrative interpretation. As these systems move into tutoring, counseling support, and creative-assistance workflows, it becomes more important to know not only whether a model can produce plausible language, but also whether its emotional judgments are calibrated against human readers. Existing multilingual profiling studies can show that language choice changes model outputs, but they do not by themselves show whether the resulting outputs are close to human judgments. Conversely, smaller human-versus-AI comparisons often focus on a single model or a single language condition and therefore do not reveal which contemporary models are actually most human-aligned \cite{shirahama2025vas,paech2024eqbench,huang2024emotionbench}.

This paper addresses that gap using a compact human-grounded benchmark for literary emotion evaluation. The core idea is straightforward. We first aggregate human Visual Analogue Scale (VAS) judgments on three short literary texts. We then compare multilingual model profiles against that human reference on four emotion dimensions: interest, surprise, sadness, and anger. The Japanese condition is treated as the primary endpoint because the human reference was collected in Japanese. English and Chinese are analyzed only as robustness conditions that test cross-language drift.

The contribution is therefore not another language-effect study in isolation. Instead, we provide a compact evaluation framework that identifies which models are closest to human reader judgments and which models maintain that closeness across languages under a common structured-scoring protocol. This framing is deliberately narrower than our prior LLM-only mathematical characterization work, which emphasized temperature-controlled and persona-conditioned individuality across 36 models \cite{shirahama2026mdpi}. Here, the priority is human alignment under a neutral-reader condition.

The present paper makes three concrete contributions. First, it provides a reproducible human-grounded benchmark design that keeps the primary claim tightly anchored to Japanese human reference data. Second, it separates primary human alignment from multilingual robustness so that secondary-language behavior is not over-interpreted as a substitute for human validity. Third, it offers a two-axis screening view---Japanese calibration versus cross-language drift---that is practical for compact engineering evaluation before larger multilingual or multi-persona studies are undertaken.
